%% ============================================================
%%  chapters/ch1_introduction.tex
%%
%%  This chapter is the EXAMPLE / TEMPLATE chapter that shows
%%  how to cite references from Zotero and how to include and
%%  reference figures.  Read every comment block carefully.
%% ============================================================

\chapter{Introduction}
\label{ch:introduction}

%% ── HOW CITATIONS WORK ──────────────────────────────────────────────────────
%%
%%  1.  Export your Zotero library (or a collection) as "Better BibTeX"
%%      (.bib format) to the file  references.bib  next to main.tex.
%%      Keep "Keep updated" checked so it auto-syncs as you add sources.
%%
%%  2.  Every source gets a unique citation key, e.g.  watsonMolecularStructure1953
%%      (Zotero's Better BibTeX generates these automatically).
%%
%%  3.  Cite in your text with  \cite{key}  for a numbered reference, or
%%      with \textcite{key}  to include the author name inline.
%%
%%  Example citation commands:
%%    \cite{key}            →  [1]
%%    \cite{key1, key2}     →  [1, 2]   (multiple sources at once)
%%    \textcite{key}        →  Watson and Crick [1]
%%    \parencite{key}       →  (Watson and Crick, 1953)  ← author-year styles
%%
%% ────────────────────────────────────────────────────────────────────────────


% ── 1.1  The DNA Molecule ───────────────────────────────────────────────────
\section{The DNA Molecule}
\label{sec:dna-molecule}

Deoxyribonucleic acid (DNA) carries the genetic information of all living organisms
\cite{watsonMolecularStructure1953}.  % ← replace key with your actual Zotero key
The double-helix structure consists of two antiparallel strands held together by
Watson-Crick base pairing: adenine pairs with thymine, and guanine pairs with
cytosine \cite{watsonMolecularStructure1953}.

Multiple references can be cited together at a single point in the text
\cite{keyOne, keyTwo, keyThree}.

You can also embed the author name naturally in a sentence using \textcite{keyOne},
who first described the phenomenon of \ldots


%% ── HOW FIGURES WORK ────────────────────────────────────────────────────────
%%
%%  1.  Place all image files (PNG, PDF, JPG, EPS) in the  figures/  folder.
%%      main.tex already sets  \graphicspath{{figures/}}  so you only need
%%      the filename, not the full path.
%%
%%  2.  Use the figure environment below.  Key commands:
%%        \includegraphics[width=…]{filename}   – import the image
%%        \caption{…}                            – figure caption
%%        \label{fig:your-label}                 – label for cross-referencing
%%
%%  3.  Refer to a figure anywhere in the text with  \ref{fig:your-label}
%%      or the clickable  \autoref{fig:your-label}  (prints "Figure 1").
%%
%%  4.  The [htbp] placement options mean LaTeX tries to place the figure
%%      Here, Top, Bottom, or on a separate Page.  Use [H] (capital) to
%%      force the figure exactly where it appears in the source (requires
%%      the float package, already loaded in main.tex).
%%
%% ────────────────────────────────────────────────────────────────────────────


% ── 1.2  Structural DNA Nanotechnology ─────────────────────────────────────
\section{Structural DNA Nanotechnology}
\label{sec:structural-dna-nano}

The predictability of Watson-Crick base pairing allows DNA to be used as a
construction material at the nanoscale \cite{seeman1982}.  Below is an example
of how to include a figure.

%% ── SINGLE FIGURE ───────────────────────────────────────────────────────────
\begin{figure}[htbp]
  \centering
  \includegraphics[width=0.7\textwidth]{example-figure.png}  % put file in figures/
  \caption{%
    A brief but descriptive caption goes here.  Always explain what the reader
    is looking at, define any labels or colours used in the image, and cite
    the original source if the figure is adapted from elsewhere
    \cite{sourceForThisFigure}.
  }
  \label{fig:dna-structure}
\end{figure}

As shown in \autoref{fig:dna-structure}, the structure demonstrates \ldots
% \ref{fig:dna-structure} also works: "Figure 1.1"


% ── 1.2.1  Holliday Junction Crossovers ────────────────────────────────────
\subsection{Holliday Junction Crossovers}
\label{subsec:holliday-junction}

The immobile Holliday junction is the key structural motif that enables the
construction of complex DNA nanostructures \cite{hollidayKey}.

%% ── SIDE-BY-SIDE FIGURES (subfigure) ────────────────────────────────────────
\begin{figure}[htbp]
  \centering
  \begin{subfigure}[b]{0.45\textwidth}
    \centering
    \includegraphics[width=\textwidth]{example-figure-a.png}
    \caption{Left panel: description of panel (a).}
    \label{fig:holliday-a}
  \end{subfigure}
  \hfill
  \begin{subfigure}[b]{0.45\textwidth}
    \centering
    \includegraphics[width=\textwidth]{example-figure-b.png}
    \caption{Right panel: description of panel (b).}
    \label{fig:holliday-b}
  \end{subfigure}
  \caption{%
    Overall caption for the two-panel figure.  Panel~(a) shows \ldots\
    Panel~(b) shows \ldots\ Adapted from \cite{sourceForThisFigure}.
  }
  \label{fig:holliday-junction}
\end{figure}

\autoref{fig:holliday-junction} shows both panels together; you can refer to
the individual panels as \autoref{fig:holliday-a} and \autoref{fig:holliday-b}.


% ── 1.2.2  2-Dimensional DNA Origami ───────────────────────────────────────
\subsection{2-Dimensional DNA Origami}
\label{subsec:2d-dna-origami}

Rothemund first demonstrated that a long scaffold strand could be folded into
arbitrary two-dimensional shapes using short ``staple'' strands
\cite{rothemundFolding2006}.

%% ── FIGURE FORCED TO EXACT POSITION ────────────────────────────────────────
%% Use [H] only when you need the figure right here, not floating.
\begin{figure}[H]
  \centering
  \includegraphics[width=0.8\textwidth]{example-origami.png}
  \caption{%
    DNA origami design showing the scaffold routing and staple positions.
    Scaffold is shown in grey; individual staples are coloured.
    Adapted from \cite{rothemundFolding2006}.
  }
  \label{fig:2d-origami}
\end{figure}


% ── 1.2.3  Three-Dimensional DNA Origami ───────────────────────────────────
\subsection{Three-Dimensional DNA Origami}
\label{subsec:3d-dna-origami}

Three-dimensional DNA origami structures are constructed by stacking multiple
2D layers \cite{douglasRapidPrototyping2009}.


% ── 1.3  Protein Structure Determination ───────────────────────────────────
\section{Protein Structure Determination}
\label{sec:protein-structure}

Membrane proteins constitute roughly 30\% of all proteins encoded in the human
genome \cite{proteinGenomeKey}.  Determining their structures is critical for
drug development \cite{drugDevKey}.


% ── 1.4  Computer Simulation ───────────────────────────────────────────────
\section{Computer Simulation}
\label{sec:computer-simulation}

\subsection{CanDo}
\label{subsec:cando}

CanDo (Computer-Aided Design of DNA Origami) \cite{candoKey} converts a
caDNAno design file into a finite-element model \ldots

\subsection{oxDNA}
\label{subsec:oxdna}

oxDNA is a coarse-grained molecular dynamics model of DNA \cite{oxdnaKey}.
Simulations were run at physiological salt concentrations \ldots


% ── 1.5  Thesis Objective ──────────────────────────────────────────────────
\section{Thesis Objective}
\label{sec:thesis-objective}

The goal of this thesis is to design, simulate, and experimentally test a DNA
origami support structure for membrane protein imaging via cryo-electron
microscopy.  The specific aims are:
\begin{enumerate}
  \item Design a half-barrel DNA origami structure using caDNAno.
  \item Simulate the design with CanDo and oxDNA.
  \item Assemble and characterize the structure experimentally.
\end{enumerate}
